\chapter{Application web et mémoire conversationnelle}

\section{API}

L'application repose sur FastAPI. L'API expose les opérations d'inscription, de connexion, de création et de consultation des conversations ainsi que l'envoi d'une question à l'agent. Elle partage le driver Neo4j et l'agent afin d'éviter plusieurs instances incohérentes. Le fournisseur de génération est transmis explicitement avec la question.

\section{Authentification}

Les utilisateurs sont enregistrés dans SQLite. Les mots de passe ne sont pas stockés en clair : une dérivation PBKDF2-HMAC-SHA256 avec sel et 600\,000 itérations est utilisée. Les sessions sont associées à des cookies. Chaque route contrôlant une conversation vérifie son propriétaire.

Les clés Mistral, Gemini et Groq ainsi que les identifiants Neo4j sont chargés depuis \code{.env}, exclu de Git. Le fichier \code{.env.example} documente les variables attendues sans contenir de secret.

\section{Mémoire}

Le graphe de connaissances et la mémoire conversationnelle remplissent des fonctions différentes. Neo4j contient les connaissances Arianes ; SQLite conserve les échanges propres à un utilisateur. Les derniers messages sont transmis à l'agent afin de résoudre les reprises telles que « et ses attributs ? ».

\section{Interface}

L'interface permet de se connecter, de créer une conversation, de sélectionner une conversation existante, de choisir le fournisseur de génération, d'afficher le schéma associé à une réponse et de dialoguer avec l'agent. Elle communique avec l'API sans accéder directement à Neo4j ou aux fournisseurs.

\capture{interface-connexion.png}{Écran d'inscription ou de connexion}{fig:connexion}
\capture{interface-conversations.png}{Interface principale et historique des conversations}{fig:conversations}
\capture{memoire-conversation.png}{Question de suivi utilisant la mémoire de la conversation}{fig:memoire}

\section{Gestion des erreurs}

Les indisponibilités de Neo4j et des fournisseurs sont distinguées. Une erreur 429 de quota, une erreur 413 de contexte ou une erreur réseau ne doit pas être présentée comme une absence du graphe. Pour une démonstration, il reste prudent de conserver des captures des réponses validées en cas de problème de proxy ou de réseau.
